%% Согласно ГОСТ Р 7.0.11-2011:
%% 5.3.3 В заключении диссертации излагают итоги выполненного исследования, рекомендации, перспективы дальнейшей разработки темы.
%% 9.2.3 В заключении автореферата диссертации излагают итоги данного исследования, рекомендации и перспективы дальнейшей разработки темы.

Во-первых, выполнен анализ предметной области автономных мобильных роботов с поддержкой периферийных вычислений. Показано, что существующие подходы к распределению ресурсов в периферийных системах развиваются по нескольким направлениям: детерминированная оптимизация, аукционные механизмы, обучение с подкреплением, многоагентное обучение и гибридные схемы. Детерминированные методы позволяют строго задавать ограничения и целевые функции, но требуют достаточно полной и устойчивой информации о состоянии системы. Аукционные механизмы позволяют учитывать дефицит ресурсов и внешние эффекты распределения, но в классическом виде ориентированы на отдельный раунд. Методы обучения с подкреплением позволяют учитывать долгосрочные последствия решений, но сами по себе не задают экономически интерпретируемого механизма распределения дефицитных ресурсов.

Во-вторых, разработана доменная и математическая модель AMR-системы с поддержкой периферийных вычислений. В модели определены множества автономных мобильных роботов, периферийных вычислительных узлов, вычислительных задач, граф сетевой связности, параметры задержек, очередей, отказов и ресурсных ограничений. Отдельно учтены классы робототехнических задач, различающиеся по критичности, срокам обработки, объёму данных и требованиям к ресурсам. Предложенная модель позволяет описывать не только вычислительную сложность задачи, но и её прикладную значимость для робота и текущей миссии.

В-третьих, построен VCG-подобный аукционный слой распределения вычислительных задач. В рамках одного раунда распределения задача назначения формализована как задача максимизации общественного благосостояния с учётом полезности выполнения задач, стоимости использования ресурсов, сетевых задержек, очередей, отказов и ограничений по срокам обработки. Платежи аукционного слоя интерпретируются как оценки внешнего эффекта, создаваемого участниками при использовании дефицитных ресурсов. При этом в работе явно зафиксированы границы применимости VCG-подобных свойств: эффективность, индивидуальная рациональность и стратегическая совместимость рассматриваются только для отдельного статического раунда при стандартных предпосылках и не переносятся автоматически на всю повторяющуюся систему с обучением.

В-четвёртых, задача адаптивного управления периферийными узлами формализована как децентрализованный частично наблюдаемый марковский процесс принятия решений. В предложенной постановке обучаемыми агентами являются периферийные вычислительные узлы, а не роботы, формирующие заявки. Действиями агентов являются операционные режимы узлов: режимы допуска, резервирования, приоритизации и защиты от перегрузки. Такое разделение ролей принципиально важно: обучаемая политика не управляет ставками и не изменяет оценки ценности задач, а значит, не подменяет аукционный механизм стратегическим обучением заявок.

В-пятых, обосновано применение кооперативного многоагентного обучения с подкреплением в парадигме централизованного обучения и децентрализованного исполнения. В качестве базового алгоритма выбран QMIX, поскольку он согласуется с дискретным пространством операционных режимов, командным вознаграждением, частичной наблюдаемостью и необходимостью локального исполнения политики периферийными узлами. Централизованное обучение может выполняться в серверном контуре, облаке, цифровом двойнике или имитационной среде, тогда как в режиме эксплуатации каждый узел выбирает режим на основе локального наблюдения.

В-шестых, разработан гибридный метод AMPLE-AMR. Метод объединяет обучаемую многоагентную политику выбора операционных режимов периферийных узлов и VCG-подобный аукционный слой распределения задач. В AMPLE-AMR обучаемая политика определяет условия распределения, то есть экспонируемую ёмкость и режимы обслуживания узлов, а аукционный слой выполняет назначение задач в рамках выбранных ресурсных ограничений. Такое построение позволяет совместить адаптацию к динамике среды и экономически интерпретируемое распределение дефицитных ресурсов.

В-седьмых, предложена масштабируемая версия C-AMPLE-AMR, основанная на кластеризации графа периферийной инфраструктуры и проведении локальных распределений внутри кластеров. Кластерная версия предназначена для случаев, когда глобальный аукционный раунд становится слишком дорогим по времени вычислений или объёму коммуникации. В работе введены критерии оценки целесообразности перехода к кластерному режиму, учитывающие накладные расходы глобального распределения, длительность управляющего слота, межкластерный срез графа и возможную потерю общественного благосостояния.

В-восьмых, реализован программный прототип AMPLE-AMR и сформирована воспроизводимая экспериментальная среда. Прототип включает генератор складских сценариев, модели роботов и периферийных узлов, эвристическое распределение, VCG-подобное распределение, кластерное VCG-подобное распределение, контур обучения режимов узлов и автоматический экспорт результатов в виде таблиц и графиков.

В-девятых, проведена экспериментальная проверка предложенного метода в сценариях стабильной нагрузки, пикового потока задач, разнородной периферийной инфраструктуры, деградации связи, отказов узлов, масштабирования и анализа чувствительности параметров режимов. Эксперименты показали, что наиболее устойчивый практический вклад вносит обучаемый выбор операционных режимов периферийных узлов. AMPLE-AMR превосходит Fixed-Heuristic по общественному благосостоянию во всех основных сценариях: на $124.0$ условных единицы при стабильной нагрузке, на $210.5$ при пиковой нагрузке, на $95.6$ в разнородной инфраструктуре, на $94.9$ при деградации связи и на $110.7$ при отказах узлов. При этом самостоятельное преимущество VCG-подобного распределения над жадной эвристикой при фиксированных режимах в текущей параметризации не подтверждено, а диагностическая роль внутренних цен подтверждена лишь частично.
